\begin{abstract}
We re-evaluate a single-cell reporting rule whose 50-cell cutpoint was declared
validated from point-estimate recovery. The original memo and calibration
criteria are reproduced in the paper. Direct evaluation on two colorectal
single-cell cohorts yields no qualifying cutpoint under pooled-tissue draws;
reference-only candidates change with cohort, count grid and aggregation.
At 50 cells, discrimination is 75.5--81.0\% in one cohort and 45.0--55.0\% in
the other, against an 80\% requirement. Controlled reuse of the same simulation
draws separates grid effects from random-stream changes. The interval target
is a draw-dependent benchmark, so these results test the recorded calibration
specification rather than population coverage. A targeted diagnostic finds 37.5--67.0\% null rejection at five cells,
explaining why high low-count exclusion rates cannot be read as reliable power.
Biological controls also fail
to establish the intended mechanistic separation. We explain the recovery
curve using an elementary estimator identity and use established trial-estimator
relationships as implementation checks. The contribution is a reproducible
failed-calibration case study: threshold validation requires direct evaluation
of intervals, null rejection and abstention, with the target and sampling
unit declared in advance.
\end{abstract}
